\section{Module 7: Introduction}

\begin{frame}
\vfill
\centering
{\color{RLBlue}\Large\bfseries Module 7}

\vspace{0.6em}
{\Large Introduction}
\vfill
\end{frame}

\begin{frame}{Why Agents, and Why Last}
\begin{itemize}
  \item By Module 3 you can build and prompt a language model; by Module 6 you can serve one efficiently. Both assume \emph{you} decide what happens next.
  \pause
  \item An \textbf{agent} is what you get when the model decides instead: which tool to call, whether to call another, and when to stop.
  \pause
  \item That single change moves every hard problem. Correctness becomes a property of a whole \emph{trajectory}, not of one output. Cost is unbounded until the loop halts. A wrong step is not discarded -- it becomes the input to the next step.
  \pause
  \item Module 7 comes last because it composes the rest: retrieval from Module 3, serving and caching from Module 6, and the evaluation discipline from Module 0.
\end{itemize}
\end{frame}

\begin{frame}{What the Eight Lectures Cover}
\begin{itemize}
  \item \textbf{The loop and its interface} (L1--L2) -- what an agent loop is, and how a model emits a call that your code can actually execute.
  \pause
  \item \textbf{What it remembers, who it delegates to} (L3--L4) -- context limits and external memory; the orchestrator--subagent pattern.
  \pause
  \item \textbf{What it retrieves} (L5--L7) -- retrieval-augmented generation end to end, embedding models, and the index structures that make search fast.
  \pause
  \item \textbf{Whether any of it works} (L8) -- trajectory-level evaluation, prompt injection, sandboxing.
\end{itemize}
\pause
\begin{rlintuition}
The first seven lectures build a system. The eighth is the only one that tells you whether it works -- which is why the project is graded on it.
\end{rlintuition}
\end{frame}

\begin{frame}{The Capstone: A Grounded Legal Research Agent}
\begin{itemize}
  \item One team project replaces the lab sequence: a research agent over national legislation that must \textbf{ground every claim in a cited article}.
  \pause
  \item Deliberately \emph{not} ``an agent that answers legal questions''. ``Find and cite passages in this corpus'' can be scored against article IDs, and it stays clear of giving advice.
  \pause
  \item \textbf{Corpus:} article-structured Uzbek legislation from \texttt{lex.uz} -- the Constitution, 20 codes and 4 major laws, $7{,}368$ articles in all.
  \pause
  \item \textbf{Why this domain:} Magesh et al.\ (2024) measured $17$--$33\%$ hallucination in commercial legal research tools marketed as eliminating it. You are graded on the harness that finds \emph{your own} rate.
\end{itemize}
\end{frame}

\begin{frame}{How the Capstone Maps onto the Lectures}
\begin{itemize}
  \item \textbf{Sprint 1} (after L1--L2) -- the agent loop, tool schemas, structured cited output. A naive baseline plus a 20-query smoke evaluation.
  \pause
  \item \textbf{Sprint 2} (after L5--L7) -- the retrieval stack: chunking, embeddings, an approximate-nearest-neighbour index, hybrid and reranked retrieval. Scored as recall@$k$ against that baseline.
  \pause
  \item \textbf{Sprint 3} (after L3--L4) -- multi-turn memory, and a verifier subagent that checks each claim against the article it cites.
  \pause
  \item \textbf{Sprint 4} (after L8) -- the evaluation harness, abstention behaviour, and a red team.
\end{itemize}
\pause
\begin{rlintuition}
Everything runs against a cheap hosted API, not a model on your laptop: cents per evaluation run, a few tens of dollars per team for the term.
\end{rlintuition}
\end{frame}

\begin{frame}{Reading Resources}
\begin{itemize}
  \item There is no textbook -- the field is younger than the publication cycle. Each lecture instead has three to five primary papers, listed on its closing slide.
  \pause
  \item \textbf{The spine:} ReAct (Yao et al., 2023) for the loop; RAG (Lewis et al., 2020) for retrieval; HNSW (Malkov \& Yashunin, 2018) for the index that makes it fast.
  \pause
  \item \textbf{Read before Lecture 8:} Magesh, Surani, Dahl, Suzgun, Manning \& Ho (2024), \textit{Hallucination-Free? Assessing the Reliability of Leading AI Legal Research Tools} -- the empirical result this module is built around.
  \pause
  \item \textbf{Benchmarks} you will meet: AgentBench, WebArena and SWE-bench for agents; LegalBench and LegalBench-RAG for the legal side.
\end{itemize}
\end{frame}
